\section{Fermat: the quotient in the manuscript and the Selmer calculation}
\label{sec:E}
We use the homogeneous plus convention
\[
 X_n:\ X^n+Y^n+Z^n=0,\qquad U_n=X_n\setminus\{XYZ=0\}.
\]
The morphism $[X:Y:Z]\mapsto[X^{n/m}:Y^{n/m}:Z^{n/m}]$
is a finite étale morphism $U_n\to U_m$ for $m\mid n$.
Its extension to the compactifications is ramified only above the
boundary. For odd $n$, $[X:Y:Z]\mapsto[X:Y:-Z]$ is the
$\mathbb Q$-isomorphism to $X^n+Y^n=Z^n$. For even $n$,
the same formula has target $X^n+Y^n+Z^n=0$ again; the
isomorphism to the minus convention requires a scalar
$\alpha$ with $\alpha^n=-1$ and the map $Z\mapsto\alpha Z$.
We keep these specified base fields throughout.

\subsection{E(a): extensions (371)--(379)}

\begin{proposition}[E.1 --- Proved: the three comparison maps]
\label{E-three-maps}
Put $B=\mathbb P^1_{\mathbb Q}\setminus\{0,1,\infty\}$,
$F=\pi_1(B_{\overline{\mathbb Q}})$, and
$N_n=\pi_1(U_{n,\overline{\mathbb Q}})$, with compatible
geometric base points for the cover
$t=-X^n/Z^n$. Then
\[
 N_n=\ker\bigl(F\longrightarrow(\mathbb Z/n)^2\bigr),
 \qquad F'=\overline{[F,F]}\subset N_n.
\]
The extensions (371), (372), and the top row of (373), in
\S51.16, folios 733v--735, scans 172--175, are respectively
\begin{align*}
 1&\longrightarrow N_n\longrightarrow\pi_1(U_n)
                 \longrightarrow G_{\mathbb Q}\longrightarrow1,\\
 1&\longrightarrow N_n\longrightarrow\pi_1([U_n/S_3])
                 \longrightarrow G_{\mathbb Q}\times S_3\longrightarrow1,\\
 1&\longrightarrow N_n/F'\longrightarrow\pi_1([U_n/S_3])/F'
                 \longrightarrow G_{\mathbb Q}\times S_3\longrightarrow1.
\end{align*}
Here $S_3$ permutes homogeneous coordinates, and the second
quotient is the fundamental group of $BS_3\times\operatorname{Spec}\mathbb Q$.
The geometric kernel in the last row is
\[
 N_n/F'\simeq n\widehat{\mathbb Z}(1)^2.
\]
The intrinsic geometric abelianization and metabelian quotient
are connected to this kernel by the continuous epimorphisms
\[
 N_n/N_n''\longrightarrow N_n/N_n'\longrightarrow N_n/F'.
\]
The middle group has rank $n^2+1$ as a
$\widehat{\mathbb Z}$-module. Thus the precise quotient in
(373) is the rank-two toric quotient. This corrects the
identification of that row in Lead 4.
\end{proposition}
\begin{proof}
Over $\overline{\mathbb Q}$ the cover adjoins $n$th roots of
$t$ and $1-t$, up to the fixed constants $-1$.
Its deck group is $(\mathbb Z/n)^2$ and its monodromy is
the reduction of $F^{\mathrm{ab}}=\widehat{\mathbb Z}(1)^2$.
This proves the kernel formula and the last kernel calculation.
Riemann--Hurwitz for a degree $n^2$ cover with three sets of
$n$ boundary points gives
\[
 2-2g(X_n)-3n=-n^2,\qquad
 g(X_n)=\frac{(n-1)(n-2)}2.
\]
The punctured surface is free of rank
$2g(X_n)+3n-1=n^2+1$. Finite étale covering theory gives
the first row. Applying it to the $S_3$-torsor
$U_n\to[U_n/S_3]$ and then to the structure map to
$BS_3$ gives the second row; it records the permutation
as well as the arithmetic action. The subgroup $F'$
is preserved by both actions, so quotienting gives
the third row. Finally $N_n''\subset N_n'\subset F'$
proves the two displayed epimorphisms. They state
exactly which information is removed, without identifying
their kernels. This checks LM7 E.13 and E.16 at the
steps used here.
\end{proof}

\begin{proposition}[E.2 --- Known, with checked reduction: sections of the minimum quotient]
\label{E-toric}
Let $T=\mathbb G_m^3/\mathbb G_m$, with coordinates
$u=X/Z$, $v=Y/Z$. The pushout of (371) by
$N_n\to N_n/F'$ is isomorphic, as an extension over
$G_{\mathbb Q}$, to $\pi_1(T)\to G_{\mathbb Q}$.
Its section classes form
\[
 H^1(\mathbb Q,\widehat{\mathbb Z}(1)^2)
       =(\widehat{\mathbb Q^\times})^2,
\]
where
\[
 \widehat{\mathbb Q^\times}
 =\{\pm1\}\times
 \left\{(a_p)_p\in\prod_p\widehat{\mathbb Z}:
 \ \text{for every }m,\ a_p\in m\widehat{\mathbb Z}
       \text{ for all but finitely many }p\right\}.
\]
The map to the base tripod's toric quotient sends a class
$(a,b)$ to $(-a^n,-b^n)$ in multiplicative Kummer notation.
For each $n$ the restriction to actual rational points of
$U_n$ is injective. The image contains formal section classes
independently of the equation defining $U_n$.
\end{proposition}
\begin{proof}
The inclusion $U_n\hookrightarrow T$ induces on geometric
groups the epimorphism sending small loops to the valuations
of $u,v$. After the identification of $N_n/F'$ with
$n\widehat{\mathbb Z}^2$, division by $n$ identifies this
map with the identity of $\widehat{\mathbb Z}^2$.
Consequently its induced morphism of quotient extensions
is an isomorphism: it is the identity on the arithmetic
quotient and an isomorphism on kernels. The identity
of $T(\mathbb Q)$ splits the target extension.
Kummer theory gives
$H^1(\mathbb Q,\mu_m)=\mathbb Q^\times/\mathbb Q^{\times m}$;
the finite groups $H^0(\mathbb Q,\mu_m)$ have the
Mittag--Leffler property, so the inverse limit has no
extra $\varprojlim^1$ term. Prime factorization gives the
displayed restricted product, since every class modulo
$m$ has finite support in its valuations.
The coordinate morphism to the base torus is
$(u,v)\mapsto(-u^n,-v^n)$, which proves the arithmetic
translation as well as the multiplication on kernels.
The map $\mathbb Q^\times\to\widehat{\mathbb Q^\times}$
is injective by valuations and sign; hence the two
coordinates recover every rational point. This is the
complete calculation of (377)--(379), scans 176--178.
It independently checks the proof of LM7 E.16.
\end{proof}

\begin{proposition}[E.3 --- Proved: equivariant splitting and its ambiguity]
\label{E-equivariant}
Write $M=\widehat{\mathbb Z}^{\,3}/\widehat{\mathbb Z}(1,1,1)$,
with the permutation action of $S_3$. The arithmetic action
on its Tate twist is scalar $\chi$.
The extension in (373) has a splitting over
$G_{\mathbb Q}\times S_3$, and any two such splittings
are conjugate by a unique element of its kernel.
Compatible choices identify its level-$n$ kernels with $nM(1)$
as in (375)--(376).
\end{proposition}
\begin{proof}
We give the cohomological verification implicit in scan 175.
The torus $T$ has identity fixed by $S_3$. Its identity
section, together with the action of $S_3$, is a section
of $\pi_1([T/S_3])\to G_{\mathbb Q}\times S_3$.
The equivariant inclusion $U_n\to T$ induces an isomorphism
from the quotient of (372) to this extension, after
identifying $nM(1)$ with $M(1)$ by division by $n$.
This proves existence at every level.

For uniqueness, $M^{S_3}=0$ and $H^1(S_3,M)=0$.
Indeed in
$0\to\widehat{\mathbb Z}\to\widehat{\mathbb Z}^{\,3}
\to M\to0$, invariants give the isomorphism between
the two diagonal copies of $\widehat{\mathbb Z}$;
$H^1$ of either permutation module is zero; and the
restriction
$H^2(S_3,\widehat{\mathbb Z})\to
 H^2(S_2,\widehat{\mathbb Z})$
is an isomorphism $\mathbb Z/2\to\mathbb Z/2$.
The latter assertion follows from
$H^2(G,\widehat{\mathbb Z})
=\operatorname{Hom}(G,\mathbb Q/\mathbb Z)$
for a finite group $G$, with the finite primary
components understood. The inflation--restriction
sequence now gives
$H^1(G_{\mathbb Q}\times S_3,M(1))=0$.
The invariant module of the product group is zero,
so the conjugating element is unique.
The same calculation holds for $nM(1)$.
For a divisibility tower the unique conjugacies commute
with transition maps. Equivalently the coordinate maps
between the tori are $n/m$ power maps and preserve the
identity. This supplies compatible splittings and
proves exactly the ambiguity asserted on scan 175.
\end{proof}

\begin{proposition}[E.4 --- Proved: the implication proposed on scan 173]
\label{E-fermat-implication}
For odd $n\geq3$ the rational boundary of $X_n$ consists
of $[0:1:-1]$, $[1:0:-1]$, and $[1:-1:0]$.
Their full cuspidal packets are disjoint from all point
sections. Therefore a proof that all sections of (371)
belong to these packets proves $U_n(\mathbb Q)=\varnothing$.
No surjectivity assertion for point sections is needed
for that implication.
\end{proposition}
\begin{proof}
A rational $n$th root of $-1$ is $-1$ for odd $n$;
setting one coordinate equal to zero gives exactly
the three displayed points. A rational point of $U_n$
always induces a section by functoriality. Stix's
Theorem 14, with the cusp decomposition embeddings,
separates this section from every packet; the
hyperbolicity condition holds already for $n=1$.
Thus a point would contradict the stated description
of the full section set. This is a proof of the
logical implication in \S51.16, folio 734, scan 173,
not an invocation of the section conjecture.
\end{proof}

\subsection{E(b): intrinsic quotients and the remaining tower problem}

\begin{proposition}[E.5 --- Known, with checked proof: the even plus curves]
\label{E-even}
For every positive even $n$, the intrinsic geometrically
abelianized extension of $U_n/\mathbb Q$ has no section.
The same holds for every quotient mapping onto its
geometric pro-$2$ abelianization, including the full
metabelian and full profinite extensions. The first
intrinsic nilpotent obstruction is at class one.
More precisely, for complex conjugation $c$,
\[
 H_1(U_n(\mathbb C),\mathbb Z_2)
 \simeq\mathbb Z_2\oplus
       \mathbb Z_2[C_2]^{\,n^2/2}
 \quad\text{as a }\mathbb Z_2[C_2]\text{-module},
\]
and the real extension represents the nonzero element
of $H^2(C_2,H_1)=\mathbb Z/2$.
\end{proposition}
\begin{proof}
We checked LM7 E.18 and E.25 and retain their actual
Schreier calculation. Positivity gives $X_n(\mathbb R)=\varnothing$.
The free antiholomorphic action on the punctured complex
surface has a nonorientable punctured quotient of Euler
characteristic $-n^2/2$. Its fundamental group $J$ is
free of rank $1+n^2/2$. Its orientation kernel $N$
has index two and rank $n^2+1$. Choose a free basis
$t,a_1,\dots,a_{n^2/2}$ for which only $t$ maps
nontrivially to $C_2$. Then
\[
 t^2,\quad a_i,\quad t a_i t^{-1}
\]
is a Schreier basis of $N$. Conjugation by $t$ fixes
the first abelian generator and swaps each remaining
pair. The square of the chosen lift of $c$ is the
first generator. Replacing $t$ by $bt$ changes this
square in the abelian kernel by $(1+c)[b]$.
Its first coordinate is always $1+2b_0$, which cannot
vanish in $\mathbb Z_2$. On a swapped pair the norm
surjects onto the invariant diagonal, so only the
first coordinate contributes to $H^2$.
Equivariant Riemann existence identifies the completion
of this finite-index extension with $\pi_1(U_{n,\mathbb R})$.
Restriction of a rational section to $G_{\mathbb R}$
and then projection to this module gives the contradiction.
Our exact program checks the ranks for $n=2,4,6,8,10$;
the Schreier calculation proves all even $n$.

The toric pushout of Proposition~\ref{E-toric} sends
this obstruction to zero: $c$ acts by $-1$ on the
torsion-free toric module, whose invariants are zero.
This explicitly relates the split extension (373)
to the obstructed intrinsic abelian extension.
\end{proof}

\subsection{E(c): the missing exponents in the Selmer theorem}

We call a section Selmer when at \emph{every} completion
its class is point-induced or cuspidal. At a real
place the point-induced class means the component
class. This is Steklov's definition in \S3.3
of~\cite{Steklov2026}; it includes $2$ and the
coefficient prime. We first prove the Selmer result,
which uses no local section conjecture.

\begin{lemma}[E.6 --- Proved: finite support forces a global point or cusp]
\label{E-finite-support}
Let $U/\mathbb Q$ be a hyperbolic curve with smooth
compactification $X$, and let $s$ be a Selmer section.
Let $Z\subset X$ be a finite reduced subscheme.
When all finite localizations of $s$ have underlying
points in $Z(\mathbb Q_v)$, the section belongs to
a point or cusp above an element of $Z(\mathbb Q)$.
In particular, a nonconstant morphism
$f:X\to W$ to a proper curve satisfying
$W(\mathbb A_{\mathbb Q})_\bullet^{\mathrm{f-cov}}=W(\mathbb Q)$
makes every Selmer section of $U$ point-induced or cuspidal.
\end{lemma}
\begin{proof}
A neighbourhood of $s$ is a finite étale cover
$U'\to U$ together with a lift of $s$.
Every lifted section is Selmer. To see the latter
directly, a point section in the subgroup corresponding
to the cover fixes a point of the finite fibre,
and so lifts to a rational point. For a cusp use the
intersection of its decomposition group with that
subgroup; its arithmetic quotient is surjective
because it contains the local section, so the
associated lifted cusp is rational.
Write $X'$ for the smooth compactification. We may
refine to genus at least one. Such refinements are
cofinal: take a characteristic geometric cover of
sufficiently large degree, spread it over a finite
extension, and take its neighbourhood twist determined
by the section. Riemann--Hurwitz, with boundary
inertia indices made arbitrarily large, gives the
required genus.

The localizations determine an adelic point of $X'$
surviving finite descent, by Harari--Stix,
Theorem 2.1. At every finite place its component
lies in the finite inverse image $Z'\subset X'$.
Stoll's Theorem 8.2 gives a rational point of $Z'$.
The system of these finite nonempty sets over all
neighbourhoods has the finite intersection property:
a common refinement maps to any finite collection
of levels. Compactness supplies a compatible
rational pro-point. Its stabilizer contains the
section; over an interior point this is the point
decomposition group, and over a boundary point it
is the cuspidal decomposition group. This proves
the first assertion.

For the second, functoriality of finite descent
sends the adelic localization to one point
$w\in W(\mathbb Q)$. The finite fibre $f^{-1}(w)$
is a suitable $Z$.
We use only finite places in this last argument.
This is the correction specified in Stoll's erratum
to Proposition 8.5; replacing a real component by
a point of the scheme-theoretic fibre would be
unjustified.
\end{proof}

\begin{proposition}[E.7 --- Proved: full level two supplies level $4\ell$]
\label{E-level-four}
Let $Z_\ell\to B$ be the finite étale cover parametrizing
a cyclic subgroup $C\subset E_\lambda[\ell]$ of the
Legendre family, for an odd prime $\ell$. There is
a nonconstant $\mathbb Q$-morphism
\[
 Z_\ell\longrightarrow Y_0(4\ell)
\]
whose extension to smooth compactifications is finite.
Every component on which an arithmetic section lifts
is geometrically connected over $\mathbb Q$.
For $\ell\geq5$ its compactification has a
nonconstant map to a positive-genus modular curve.
\end{proposition}
\begin{proof}
The ordered points $P=(0,0)$ and $Q=(1,0)$ of order
two are independent on $E=E_\lambda$.
Let $\phi:E\to E'=E/\langle P\rangle$ be the quotient
and let $\phi':E'\to E$ be its dual, so
$\phi'\phi=[2]$ and $\phi\phi'=[2]$.
Define
\[
 H_4=(\phi')^{-1}(\langle Q\rangle)\subset E'.
\]
It has order four. For a point $R$ with
$\phi'(R)=Q$, we have $2R=\phi(Q)\ne0$,
so $R$ has order four and generates $H_4$.
Thus $H_4$ is cyclic, defined over the base, and
$\phi(C)$ has order $\ell$. Their sum is a cyclic
subgroup of order $4\ell$, giving the asserted
moduli map. All constructions commute with
base change; the finite étale subgroup scheme
$C$ is the tautological one on $Z_\ell$.
The function $j(E')$ is nonconstant: a fixed elliptic
curve has only finitely many degree-two isogenous
curves, whereas $j(E_\lambda)$ is nonconstant.
Therefore the map of smooth projective curves
is nonconstant and finite.

For completeness, a connected arithmetic cover
supporting a section onto $G_{\mathbb Q}$ has
constant field $\mathbb Q$: otherwise its
fundamental group's arithmetic image is the
proper subgroup attached to the constant field.
For $4\ell$ the modular index is $6(\ell+1)$,
there are six cusps, and there are no elliptic
points of orders two or three. Consequently
\[
 g(X_0(4\ell))=1+\frac{6(\ell+1)}{12}-\frac62
             =\frac{\ell-3}{2}.
\]
At $\ell=5$ the target is $X_0(20)$ of genus one.
The exact program checks these indices and genera
at the first seven odd primes. The displayed
formula proves the general assertion.
\end{proof}

\begin{lemma}[E.8 --- Known, with checked steps: the local residual representations]
\label{E-local-representations}
Let $\ell\geq5$ be prime and let $s$ be a Selmer
section of $U_\ell/\mathbb Q$. Let
$\bar\rho_s:G_{\mathbb Q}\to\mathrm{GL}_2(\mathbb F_\ell)$
be obtained from the Legendre family.
It is reducible. There is a quadratic twist whose
semisimplification is $1\oplus\bar\chi_\ell$.
The mod-two representation is trivial before
and after that twist.
\end{lemma}
\begin{proof}
We checked Steklov's Lemmas 5.4, 5.7, 5.11--5.12 and
Propositions 5.5, 5.9--5.10, 5.13, including their
uses at $2$ and $\ell$.
Here are the arithmetic steps needed for the
extension to $\ell=5$.
For a local point write
$a^\ell+b^\ell+1=0$ and $\lambda=-a^\ell$.
The formula
\[
 j=2^8\frac{(\lambda^2-\lambda+1)^3}
                {\lambda^2(\lambda-1)^2}
\]
gives $v_p(j)=8v_p(2)-2\ell|v_p(a)|$ when
$v_p(a)\ne0$. When $v_p(a)=0$ and $\lambda\equiv1$
use $1-\lambda=-b^\ell$ instead. At $2$ one of
these cases has positive absolute valuation;
$\ell\geq5$ makes $v_2(j)<0$. At odd $p$ there
is either good reduction or potential multiplicative
reduction with $\ell\mid v_p(j)$.
After a local quadratic twist, the latter is a
Tate curve. At $p\ne\ell$ its residual extension
class is the Kummer class of $q$, so divisibility
of $v_p(q)$ by $\ell$ makes it unramified.
At $p=\ell$, the same divisibility puts $q$
modulo $\ell$th powers in the unit subgroup.
The residual representation is finite flat,
hence has Serre weight two; good reduction has
the same weight. At $2$ the untwisted Tate
representation is tame, with conductor exponent
at most two (the bound four used by Steklov
also suffices).

For a local cusp, the map to the $j$-line has
ramification index divisible by $\ell$.
The Tate parameter class is therefore divisible
by $\ell$. Its residual representation is a
quadratic twist of $1\oplus\bar\chi_\ell$.
In particular the same local conditions hold,
including at $2$ and $\ell$.

Suppose $\bar\rho_s$ were irreducible.
Its determinant is $\bar\chi_\ell$, so complex
conjugation has distinct eigenvalues $1,-1$.
Any invariant line over $\overline{\mathbb F}_\ell$
would be one of those two lines, already defined
over $\mathbb F_\ell$; hence irreducibility is
absolute. The local quadratic characters on
inertia patch to a global quadratic character:
the abelian cyclotomic isomorphism
$G_{\mathbb Q}^{\mathrm{ab}}\simeq
 \prod_p\mathbb Z_p^\times$
allows us to multiply their finitely many
nontrivial inertia characters. The twisted
representation has determinant $\bar\chi_\ell$,
weight two, and level dividing $16$.
The Khare--Wintenberger theorem then produces
a weight-two cusp form of level dividing $16$.
Each such $X_0(N)$ has genus zero, a contradiction.

Thus the representation is reducible.
At $\ell$, a good supersingular representation
over $\mathbb Q_\ell$ has irreducible inertial
action, so this case is excluded. The local
semisimplifications are quadratic twists of
$1\oplus\bar\chi_\ell$. Patch again and write
the two global characters as $\omega$ and
$\bar\chi_\ell\omega^{-1}$.
They are unramified away from $\ell$, and one
of them is unramified at $\ell$. A finite
extension of $\mathbb Q$ unramified at all
finite primes is trivial by Minkowski's
discriminant theorem. Hence one character is
$1$ and the other $\bar\chi_\ell$.
Finally the three Legendre points of order
two are rational, and a quadratic twist acts
as $[-1]$, which is the identity on $E[2]$.
\end{proof}

\begin{theorem}[E.9 --- Proved: the Selmer theorem includes exponent five]
\label{E-selmer-primes}
For every prime $\ell\geq5$, every Selmer section
of $U_\ell/\mathbb Q$ is cuspidal.
\end{theorem}
\begin{proof}
The invariant line of Lemma~\ref{E-local-representations}
lifts the pushed-forward section on $B$ to the
cover $Z_\ell$ of Proposition~\ref{E-level-four}.
The lift is Selmer by the finite-fibre argument
in Lemma~\ref{E-finite-support}.
Its compactification maps nonconstantly to
$X_0(4\ell)$, whose positive genus and Stoll's
Corollary 8.8 give
\[
 X_0(4\ell)(\mathbb A_{\mathbb Q})_\bullet^{\mathrm{f-cov}}
        =X_0(4\ell)(\mathbb Q).
\]
Lemma~\ref{E-finite-support} therefore makes the
lift, and hence the section on $B$, point-induced
or cuspidal. A point-induced section on $B$
lifts to a rational point of $U_\ell$ whenever
its image is contained in the subgroup
$\pi_1(U_\ell)$: it fixes an element of the
finite étale fibre. A cuspidal section lifts
into the intersection with a cusp decomposition
group and is cuspidal upstairs.

It remains to exclude the point case.
The associated rational Legendre elliptic
curve has full rational two-torsion.
After the quadratic twist of
Lemma~\ref{E-local-representations}, its
$\ell$-torsion has a submodule $1$ or
$\bar\chi_\ell$. In the first case it has a
rational point of order $\ell$. In the second,
quotient by the subgroup $\mu_\ell$; the
dual isogeny's kernel is the constant group
$\mathbb Z/\ell$, so the quotient elliptic
curve has a rational point of order $\ell$.
The odd-degree isogeny is an isomorphism on
two-torsion. Thus one rational elliptic curve
has a subgroup
$\mathbb Z/2\oplus\mathbb Z/(2\ell)$ in its
rational torsion. Mazur's Theorem (8)
allows such a subgroup only for $\ell\leq4$.
This contradiction proves the assertion.

The point-case exclusion also repairs the
last sentence of Steklov's Theorem 5.14.
His claim that $Y_1(\ell)(\mathbb Q)$ is empty
for $\ell\geq7$ is false at seven.
The necessary extra datum is the full
two-torsion, retained throughout this proof.
\end{proof}

\begin{proposition}[E.10 --- Refuted, with repair: the claimed emptiness at seven]
\label{E-seven-error}
The claim $Y_1(7)(\mathbb Q)=\varnothing$ used in the
last sentence of Steklov's Theorem 5.14 is false.
The elliptic curve
\[
 E:\ y^2-xy-4y=x^3-4x^2
\]
has discriminant $-1664$ and the point $(0,0)$
of exact order seven. The simultaneous requirement
of full rational two-torsion and rational
seven-torsion is impossible by Mazur's Theorem (8).
\end{proposition}
\begin{proof}
The successive nonzero multiples of $(0,0)$ are
\[
 (0,0),\ (4,8),\ (2,2),\ (2,4),\ (4,0),\ (0,4);
\]
the seventh multiple is the origin.
The included exact rational-arithmetic program
checks the curve equation at every step, the
addition law, and the nonzero discriminant.
The last assertion follows from the forbidden
subgroup $\mathbb Z/2\oplus\mathbb Z/14$.
This is precisely the replacement used in
Theorem~\ref{E-selmer-primes}.
\end{proof}

\begin{theorem}[E.11 --- Proved: exponent three]
\label{E-selmer-three}
Every Selmer section of $U_3/\mathbb Q$ is cuspidal.
\end{theorem}
\begin{proof}
On $X_3$ the functions
\[
 u=-\frac{3Z}{X+Y},\qquad
 v=\frac{9(X-Y)/(X+Y)-1}{2}
\]
give a birational map to $E:v^2+v=u^3-7$.
Indeed multiplication by $(X+Y)^3$ turns
$v^2+v-u^3+7$ into
$27(X^3+Y^3+Z^3)$, as the exact symbolic
calculation verifies. The smooth projective
curves therefore are isomorphic. The map sends
the three rational boundary points to
$(3,-5)$, $(3,4)$, and the origin.

The change of variables
$(x,y)=(4u,8v+4)$ gives $y^2=x^3-432$.
The identification of this curve with $X_0(27)$
is explicit in Kriz--Li, \S10.1~\cite{KrizLi2019};
it is the curve $27a1$. Stoll's Corollary 8.8
gives its finite-descent equality and rank zero.
For the rational torsion, good reduction at
$5$ and $7$ gives respectively $6$ and $9$
points. Prime-to-residue-characteristic
torsion injects under reduction. Applying
the two reductions prime by prime bounds
the rational torsion by three; the three
displayed points give equality. Thus
$X_3(\mathbb Q)$ consists exactly of these
boundary points. A Selmer section localizes
to their diagonal image under the
finite-descent equality. Lemma~\ref{E-finite-support}
then puts it into a boundary packet.
\end{proof}

\begin{corollary}[E.12 --- Proved: all exponents and the stipulated local input]
\label{E-all-exponents}
For every odd $n\geq3$, all Selmer sections of
$U_n/\mathbb Q$ are cuspidal. For every even
$n\geq2$ there are no sections at all.
Using the local theorem stipulated as Tool (1)
in the development prompt, every section of
$U_n/\mathbb Q$ is Selmer. Consequently,
in the stipulated framework the complete
section set for odd $n\geq3$ is the disjoint
union of the three rational cusp packets,
each a torsor under
$H^1(\mathbb Q,\widehat{\mathbb Z}(1))$.
\end{corollary}
\begin{proof}
Choose an odd prime $\ell\mid n$. The finite étale
power map $U_n\to U_\ell$ carries local point
sections to point sections and local cusp
sections into cusp packets, so it preserves
the Selmer property. Its image is cuspidal
by Theorems~\ref{E-selmer-primes} and
\ref{E-selmer-three}. Intersecting its
decomposition group with $\pi_1(U_n)$ gives
a cusp decomposition group whose image in
$G_{\mathbb Q}$ is surjective because it
contains the given section. The cusp is
therefore rational and the section is
cuspidal. The even assertion is
Proposition~\ref{E-even}.

For the final assertion, Tool (1) and the
neighbourhood reduction in Part B give the
point-or-cusp description at every finite
place, with no exclusion of $2$ or $\ell$.
The real section theorem gives the point
component at the infinite place.
These are exactly all the places required
by Steklov's definition. The three packets
exist by choosing local parameters; their
disjointness and their torsor descriptions
are Stix's Theorem 14 and Proposition 8.
This proof identifies where the unrefereed
2026 local input enters. The Selmer theorem
above is independent of that input.
\end{proof}

\paragraph{Provenance of the local input.}
The final passage in Corollary~\ref{E-all-exponents}
uses the supplied October 2026 manuscript as
the tool explicitly stipulated in the question.
Our checking of its named intermediate steps
is recorded in Part B and the audit.
We do not report an independent verification
of every argument in that manuscript.
The finite-level tower problem in E(b) is
a further task: the morphism from the full
section set to a quotient's section set need
not be surjective. The next subsection keeps
its image and its obstruction groups explicit.

